% ================================================================
%  第 7 章 双重面孔 (原书 p.153-166, 边码 p.37-p.56)
%  源文: _work/src1/task-ch07.md (= ch07.md)
% ================================================================
\chapter{双重面孔}

\section{(16) 沼泽与头几排}

% ---------------- p.153 ----------------
\origpage{37}可我还没有走到这场反思的尽头:反思的是,在轻蔑的出现与蔓延之中,我占了怎样一份——在这样一个世界里,我至今仍兴高采烈地以“数学共同体”之名来指称它。我现在感觉到,正是这场反思,才是我所能献给那个世界里我所爱过的人的最好的东西;此刻我正准备着——当然不是回到那个世界去,而是要在那里重新开口说话。

我想,我尤其还剩下这样一件事要做:审视一下,在我同他们一样仍是那个世界中一员的那些年月,我与这个那个成员之间,维持的是怎样一种关系。

如今回想起来,有一点让我吃惊:那个世界里有一整块人,我虽说常与他们擦肩而过,却始终没进入我的注意,仿佛他们并不存在。当年我大概把他们感知为某种“沼泽”——在我心里并没有明确的功能,我想,连“共鸣箱”的功能都谈不上;像一团灰色的、无名的人群:在讨论班和学术会议上,他们一成不变地坐在最后几排,仿佛生来就被分派在那儿;报告进行中,他们从不开口冒险提一个问题,因为他们事先就确信无疑:自己的问题只可能问得不着边际。倘若他们向像我这种被公认为“在行”的人提问,那准是在走廊里,在“权威们”显然无意彼此交谈的时候——这时他们便飞快地、像踮着脚尖似的把问题提出来,仿佛为占用我们这种大人物的宝贵时间而羞愧。有时那问题确实显得不着边际,我便(我想)试着用几句话说明为什么;常常它也是切中要害的,我同样尽力作答,我相信。两种情形下都一样:在这样一种心态里(或者我该说,在这样一种气氛里)提出的问题,很少会有第二个问题跟上,把它说得更确切、更深入。也许我们这些坐在头几排的人,在那种场合确实太匆忙了(尽管我们有时无疑也在努力不显出匆忙的样子),以至对面那份畏惧无法消散,交谈也就无从产生。我当然感觉得到,正如我的对话者在他那一侧也感觉得到:我们卷入的这局面有多么虚假、多么造作——只是我那时从未对自己把它说出来,而他也一样,想必从未说出来过。我们两个都像古怪的自动机那样运转,一种古怪的默契把我们连在一起:假装不知道那攫住我们之中一个人的焦虑,而另一个人则隐隐约约地觉察到了它——那是充满焦虑的空气里的一小份焦虑,这空气浸透了那些场所,所有人想必都像我们一样觉察得到,又都选择一致同意地视而不见\fnoteP{1}{(13)}{前面这段描述,显然只求勉强用具体的字眼,把记忆的那团“雾”交给我的东西还原出来;那团雾从未凝结成任何一个稍微确切的具体情形,好让我在这里给出稍微“写实”或“客观”一点的描述。若让这段文字说出这样的意思——那些不愿坐在头几排、或者没有明星与“权威”身份的同行,在同这样一位人物说话时必定满心焦虑——那就是曲解我的话了。在这个圈子里我认识的大多数朋友,显然都不是这样,甚至那些偶尔出入学术会议或讨论班的人也不是。毫无保留地成立的是:“权威”身份造成一道屏障、一道沟壑,把它与没有类似身份的人隔开;而这道沟壑极少消失,哪怕只是一场谈话的工夫。我还要补充:“头几排”与“沼泽”之间那种主观的区分(在我看来却实实在在),绝不能化归为社会学标准(社会地位、职位、头衔等等),甚至也不能化归为“身份”、名声的标准;它还映照出气质上、或者更难把握的禀性上的一些心理特点。我二十岁到巴黎时,就知道自己是个数学家,我做过数学;尽管有我曾有机会谈过的那些水土不服,我心底里觉得自己是“他们中的一员”——只是这一点只有我自己知道,而且我起初甚至并不确信自己还会继续做数学。如今我倒更倾向于坐到最后一排去(在难得出现这个问题的场合)。} (13)。

% ---------------- p.154 ----------------
\origpage{38}这种对焦虑的模糊感知,在我这里要等到1970年第一次“觉醒”之后才变成自觉的;那时,这片“沼泽”从我乐意让它一直待着的那片阴影里走了出来。这件事并不是出于什么有意的决定,我也没有立刻意识到,我就这样离开了一个圈子,进入另一个圈子——从“头几排”人们的圈子进入“沼泽”:突然间,我的新朋友大多正是那些一年前我还会不声不响地归入那片无名无状的疆域的人。所谓沼泽,忽然间动了起来,借着一张张朋友的面孔活了过来——这些朋友与我由一场共同的历险连在一起,另一场历险!

\section{(17) 特里·米尔基尔(Terry Mirkil)}

说真的,早在这场关键的转折之前,我就与几位同学(后来成了“同行”)结下了友谊,假如当年有人问我,我无疑会把他们归入“沼泽”(而且假如他们不是我的朋友……)。非得有这场反思、非得我去翻检记忆,我才想起他们来,那些零散的记忆才聚拢到一处。我认识这三位朋友,是在最初的那些日子里;那时我和他们一样,在南锡学着这门手艺——也就是说,那时我们还在同一条船上,还没有任何东西把我标举为一位“权威”。这多半不是偶然;而此后二十年里再没有过别的这样的友谊,多半也不是偶然。我们四个都是外国人,这肯定是一条不可忽视的纽带——我与那些和我一样被“空投”到南锡的年轻“高师人”的关系,就远没有那么私人化;我们几乎只在系里见面。三位朋友中有一位,一两年后移居南美。他和我一样是法国国家科学研究中心(CNRS)的研究专员,而我总有这么个印象:他自己也不太清楚他在“找”什么;久而久之,他在 CNRS 的处境变得有点尴尬。我们继续偶尔见面或通信,最后失去了联系。我与另外两位朋友的关系持续得更久,也更强烈,远不那么浮浅。何况数学上的共同兴趣在其中只扮演了极黯淡的角色,甚至等于零。

% ---------------- p.155 ----------------
\origpage{39}\origpage{40}与特里·米尔基尔(Terry Mirkil)和他的妻子普雷索西亚(Presocia)——她娇小纤弱,正如他粗壮结实;两个人都带着一股温柔的神气——我们在南锡常常一晚上、有时一整夜地唱歌、弹钢琴(那时是特里弹),谈音乐(那是他们的挚爱),也谈我们各自生活里这样那样重要的事情。老实说,不是最重要的那些——不是那些总被如此小心地闭口不谈的事。然而这份友谊给了我很多。特里有一份我所缺乏的细腻与辨别力,而我大部分精力早已倾注在数学上了。远比我更多地,他保留着对单纯而本质之物的感觉——太阳、雨、土地、风、歌唱、友谊……

特里在达特茅斯学院(Dartmouth College)找到一个合他心意的职位之后——那里离哈佛不算太远,从五十年代末起我常去那儿小住——我们仍继续见面、通信。其间我得知他容易陷入抑郁,这让他一次次长时间住进“疯人院”——在他唯一一次、也是极简略地向我提到此事的信里,他这样称呼那些地方;那是在又一次“可怕的长住”之后。我们见面时,从不谈这件事——只有一两次极偶然地,为了回答我的惊讶:他和普雷索西亚为什么不收养一个孩子。我不认为我脑子里曾冒出过这样的念头:他和我,我们可以谈谈问题的底;或者哪怕只是碰一碰——恐怕连“也许有些问题值得看看”这个念头也不曾有过,无论在我朋友的生活里,还是在我自己的生活里……在这些事情上,有一个不曾说出口、也无法逾越的禁忌。

渐渐地,见面和通信稀疏了。确实,我越来越成了工作和一个角色的囚徒,尤其是成了那份意志的囚徒——它变成了一种固执的念头,也许是对别的东西的一种逃避:在作品的累积中不断地超越自己——而与此同时,我的家庭生活正神秘地、不可挽回地败坏下去……

有一天,我从特里在达特茅斯的一位同行来信中得知,我的朋友自杀了(那已是他死去并被埋葬很久之后……);这个消息像穿过一层雾传到我这里,像来自一个极其遥远、而我不知何时已经离开了的世界的回声。也许那是我心里的一个世界,它早在特里结束自己生命之前就已死去——他的生命被一种焦虑的暴力摧毁,那焦虑他没能、或者不愿去化解,而我也没有能够、或者不愿意去猜出……

\section{(18) 二十年的自命不凡,或:不知疲倦的朋友}

% ---------------- p.156 ----------------
我与特里的关系,我想,在任何时刻都不曾因我们在数学界地位的差异、或因我由此而生的优越感而变味。这份友谊,以及那些年月里生活馈赠给我的另外一两位朋友(全然不管我“配不配”),肯定是当时少有的解药之一,用来对付一种隐秘的自命不凡——它由社会地位滋养,更由我对自己数学力量的自觉以及我自己给予它的价值滋养。我与第三位朋友的关系就不同了。这位朋友,以及后来他的妻子(他是在我们于南锡相识前后认识她的),在这些年里,每一次我们相遇,无论在他们家还是在我家,都对我表现出一种温暖的友谊,满含体贴与质朴。这份友谊里显然没有任何与地位或脑力相关的潜台词。然而,我与他们的关系,二十多年里始终带着我内心那种深刻的暧昧,带着我说过的那道分裂——它给我的数学家生活打上了印记。在他们面前,每一次我都忍不住感到他们亲切的友谊并作出回应,几乎是违心的!与此同时,二十多年里我居然做到了这样一手绝活:从我的伟大高处,轻蔑地俯视我的朋友。这想必从在南锡的头几年就这样开了头;而且很长时间里,我的成见也扩展到了他妻子身上,仿佛\origpage{41}\origpage{42}事先就可以认定:他的妻子只可能同他一样“无足轻重”。母亲和我之间,我们喜欢只用一个嘲弄的绰号来指他;这绰号想必在我母亲1957年去世之后很久还刻在我心里。如今我看出,至少在我这态度背后的诸般力量之一,是我母亲那坚强的个性在她一生中、乃至她死后近二十年里对我的支配——那些年里,我继续浸染在支配了她一生的那些价值之中。我朋友那温顺、和蔼、毫无争强好胜之心的天性,被不声不响地归入“无足轻重”一类,成了嘲讽式轻蔑的对象。直到现在,我第一次费心去审视这段关系究竟是什么样子,才发现这种面对他人温暖善意而执意自我隔绝的状态,范围有多大——它曾如此长久地给这段关系打上印记。我的朋友特里,并不比另一位朋友更好斗、更犀利,却幸好得到了我母亲的认可,不曾成为她嘲笑的对象——我猜,正因如此,我与特里的关系才得以在我心里毫无阻力地绽放。他对数学的投入并不更热忱,“天赋”也并不更出众,可我并没有因此就找到借口,用那层轻蔑与自得的硬壳把自己同他和他的妻子隔开!

在这另一段关系里,我至今仍无法理解的,是我朋友那份亲切的友谊,何以从未因他每次见面都必定感觉到的我的迟疑而灰心。然而今天我很清楚:除了那层硬壳和那份轻蔑,我还是别的东西;除了一块脑肌肉和由它生出虚荣的自命不凡,我还是别的东西。和他们一样,我心里也有那个孩子——那个我假装不知道的孩子,轻蔑的对象。我把自己与他隔断了,可他仍在我心里的某处活着,健康而茁壮,像我出生那天一样。我朋友们的感情,肯定是给那个孩子的——他们不像我那样与自己的根隔断。而肯定也正是他,在大老板背过身去的时候,偷偷地、匆匆地回应了那份感情。

\section{(19) 没有爱的世界}

% ---------------- p.157 ----------------
大老板老了,幸好,他也稍稍风化了一些,那小子从此可以自在些了。至于同这些真正经得起考验的朋友的关系,我想我在这里确实指到了我一生中最触目、最怪诞的一例:某种自命不凡(除了别的之外)在一段私人关系里所起的效用。也许我还在自欺,但我相信这也是唯一的一例:我与数学界(甚至别处)某位同行或朋友的关系,被自命不凡持久地占据,而不是让它满足于偶尔地、隐约而短暂地露一露头。此外我似乎觉得,我当时在数学界乐于来往的众多朋友当中,没有一个我能想象他也经历过类似的事——在一段与同行(朋友也罢,不是也罢)的关系里。在我所有的朋友当中,我也许是最不“酷”的,最“死用功”的,最不肯让一丝幽默透出来的(这本事我晚年才终于有了),最容易把自己看得无比严肃的。肯定也正是因此,我才不会太去寻求与我同类的人作伴(假定真有这样的人)!

奇怪的是,我的朋友们,不管属于“沼泽”与否,都容忍我,甚至喜欢我。这是一件好事,也是一件要紧的事,值得在这里说出来——尽管我们常常一连几小时、几天只见面讨论数学:感情在流动,就像今天它仍在流动一样,在当时的友人(凭一些有时纯属偶然的意气相投)与我之间流动;从1949年在南锡我被亲切接纳的那一刻起——在洛朗·施瓦茨(Laurent Schwartz)与伊莱娜·施瓦茨(Hélène Schwartz)家里(我在那儿有点像家里的一口人),在让·迪厄多内(Jean Dieudonné)家里,在罗杰·戈德芒(Roger Godement)家里(有一阵我也常去\origpage{43})。

这种包裹着我在数学界最初几步的温暖感情,我有点倾向于把它忘掉,它对我整个数学生涯都很重要。肯定是它,让我同那个以长辈们为化身的圈子的关系,也带上了同样温暖的色调;它给了我对这个圈子的认同以全部力量,也给了“数学共同体”这个名号以全部意义。

显然,对今天许多年轻数学家来说,事情是这样:在他们学习的年月、往往还远不止这些年月里,他们被切断了与一切感情、温暖之流的联系;他们的工作只映在一位疏远的老板眼中,映在他吝啬的评语里,有点像在读科研与工业部的一份通令——那通令剪掉了工作的翅膀,剥夺了它比一份阴郁而不牢靠的糊口营生更深的意义。

可我说着说着就提前了——说到七十、八十年代数学界那种失宠,也许是所有失宠中最深的一种;那个数学界里,发号施令的是我昔日的学生,以及我旧日朋友的学生。在那个世界里,老板常常把工作题目分派给学生,像扔给狗一根骨头——就这个,别的没有!像给囚犯分派一间牢房:你就在那儿服你的孤独刑!在那个世界里,某项细致而扎实的工作、多年耐心努力的结果,被那位无所不知、手握权力的人微笑着轻蔑地否掉:“这工作不合我胃口!”事情就这样定了。扔进垃圾桶吧,别再提了……

这样的失宠,我很清楚,在我所认识的那个圈子里、在我五十、六十年代常来往的朋友中间,并不存在。确实, 1970年我得知:在数学之外的科学界,这倒更像是家常便饭——甚至在数学里,显然也不那么罕见:公开露骨的轻蔑,明目张胆的滥用权力(而且无处申诉),甚至在某些有名望、我也曾有机会遇见的同行身上。但在我天真地当作“那个”数学界、或者至少当作这个世界忠实缩影的那个朋友圈子里,我没有见过这样的事。

然而,轻蔑的种子想必已经在那里了,由我的朋友和我播下,在我们的学生身上发了芽。而且不仅在我们的学生身上,也在我某些旧日的伙伴和朋友身上。但我的角色不是揭发,甚至不是斗争:腐败是不能靠斗争解决的。在我爱过的某个学生身上、或在某个旧日伙伴身上看到它,我心里有什么东西会揪紧——而往往,我不是去接受痛苦带给我的认识,而是拒绝这痛苦,挣扎,躲进拒绝和一种斗争姿态里:这种事不该存在!可它偏偏存在——而且我心底知道它的意义何在。从不止一个意义上说,我与此并非无关:我所爱的某个学生或旧日伙伴,乐于不声不响地踩一踩另一个我所爱、并且他在其中认出我的人。

我又跑题了,可以说跑了两重题——仿佛轻蔑的风只绕着我的住处吹!然而,正是它吹在我身上、吹在我亲近和珍爱的人身上的那口气,使我被它触及,认识了它。但现在还不是谈这个的时候,除非只对我自己,在沉默里谈。现在该是我重新拾起这场反思—见证的线索的时候了;它不妨就叫“追踪轻蔑”——追踪我心中和我周围的轻蔑,在我曾属于的那个数学界里,在五十、六十年代。

\section{(20) 一个没有冲突的世界?}

% ---------------- p.158 ----------------
\origpage{44}\origpage{45}我原打算用几行字谈谈“沼泽”,算是尽一份心,只说它在那儿,而我不常去——就像沉思中(数学工作里也如此)常有的事:你注视的那个“什么也没有”,结果却满是生命、奥秘,以及一向被忽略的认识。就像另一个“什么也没有”,它也像碰巧似的就在南锡(看来南锡真是我这个新身份的摇篮!):那个多半有点不中用的学生的“什么也没有”——人家是怎么对待他的,那可真该瞧瞧……刚才我还闪过一下这个念头,就在我写下(也许写得太快了点?)“这些失宠”在我们这儿“还不存在”的时候。这么说吧:我能报告的这类事,只有这独一无二的一桩;它像(这总得承认)我暗指的那种“失宠”,而我不愿过多逗留在细节描述上。受过这种对待的人,用不着我画图,就很清楚我说的是什么。那些没受过、但每次遇上都不急于闭上眼睛的人,也一样清楚。至于其余的人,那些穷开心地轻蔑别人的人,以及那些满足于闭上眼睛的人(就像我自己二十年间成功地做到的那样),哪怕给他们画一本画册也是白费……

我还剩下要做的是:审视这两个十年里我与同行、与学生之间的私人与职业关系,顺便也审视一下我所知道的、我最亲近的同行彼此之间以及他们与学生之间的关系。今天最让我吃惊的是:在所有这些关系里,冲突似乎根本不存在。我得马上补充:这在当时我觉得再自然不过——差不多是最起码的事。冲突,在善良意愿的人之间、在智性与精神上都已成年的人之间、以及诸如此类的人之间(再说一遍:最起码的事!),本就不该有。哪里真有了冲突,我便把它看作一种令人遗憾的误会:只要有应有的善意,把话说清楚,那准能在最短时间内了结,不留痕迹!我年轻时之所以选择数学作我最喜爱的活动,肯定是因为我感到:在这条路上,这种世界观最有可能不至于每每一头撞上令人不安的反证。毕竟,当人们证明了一件事,所有人就一致了——也就是说,善良意愿的人以及诸如此类的人,这自然不在话下。

碰巧,我确实感觉对了。这两个十年,是在我亲爱的“数学共同体”那个“没有冲突”(?)的世界的安宁里度过的;这段历史,也是我内心长期停滞的历史:眼睛和耳朵都塞住了,除了数学几乎什么也没学到——而在我的私人生活里(先是我和母亲之间的关系,然后是我在她死后随即建立的家庭),一场无声的毁灭正猖獗横行,那些年里我任何时候都不敢去看它一眼。可那是另一个故事了……我在这几行里常提到的1970年“觉醒”,不仅是我数学生涯的转折、一次圈子的彻底更换,也是(差一年)我家庭生活的转折。也正是那一年,我第一次在与新朋友的接触中,冒险偶尔——还非常偷偷地——瞥一眼我生活里的冲突。正是那个时候,一个疑问开始在我心里冒出,并在随后几年里成熟:我生活里的冲突,以及我有时在别人生活里感到的那种冲突,并不只是一场误会、一道用海绵一抹就掉的“污痕”。

这种(至少相对的)冲突之缺席,在我选为己有的那个圈子里,事后看来是相当了不起的事——而我后来终于知道:凡有人类生活的地方,冲突都肆虐横行,在家庭里如此,在工作场所也如此,无论是工厂、实验室,还是教授或助教的办公室。简直可以说, 1948年9月或10月,我懵懵懂懂地在巴黎下船,正好一头栽进了宇宙间唯一那座天堂般的小岛,那里的人们彼此相处没有冲突!

% ---------------- p.159 ----------------
\origpage{46}一下子,我觉得这事真了不起——在1970年以来我所知道的一切之后。它当然值得仔细看看——这是神话,还是现实?我看得见流动在我众多朋友与我之间、后来又流动在学生与我之间的那份感情,我用不着去编造它——可在这个仿佛把冲突放逐了的天堂世界里,我几乎不得不去编造冲突!

确实,在这场反思里,我毕竟有机会碰到过两种冲突的情境,每一次都像是我内心某种态度的显影剂:一个是南锡“不中用的学生”那桩事,其中当事人之间的是非曲直我并不清楚。另一个是我自己内心的冲突情境,一道分裂,在我与“不知疲倦的朋友”的关系里——但这一桩从未以人与人之间冲突的形式表现出来,而那本是冲突唯一被普遍承认的形式。值得注意的是,按这个词的通常意义,这些朋友与我之间的关系完全没有任何冲突——它任何时候都不曾有过一丝阴云。分裂在我心里,不在他们那里。

我继续清点。最先想到的之一:布尔巴基团体!在我或多或少定期参加它的那些年里——也就是直到五十年代末——这个团体对我而言体现了一种集体工作的理想:既尊重这项工作中看似极细微的细节,也尊重每个成员的自由。在布尔巴基团体的朋友们中间,我任何时候都不曾感到一丝强制的气味,无论是对我,还是对任何别人,无论资深成员还是受邀而来、试试看他和团体之间是否“来电”的客人。任何时候都不曾有过一丝影响力之争,无论是围绕议程上某个问题上的观点分歧,还是围绕对团体的霸权。这个团体没有首领而运转,而且据我所能察觉,显然没有人内心深处渴望扮演这个角色。当然,如同一切团体,某个成员对团体、或对其他某些成员的影响力大于别人。韦伊在这件事上扮演了一个特殊的角色,我谈过。他在场时,他有点像“主持人”\fnoteP{2}{(14)}{人们会以为这与“没有首领”的说法相矛盾,其实不然。对布尔巴基的老成员来说,我想,韦伊被看作这个团体的灵魂,但从来不是“首领”。他在场而且他乐意的时候,就如我说过的,他成了“主持人”,但他并不发号施令。他心情不好的时候,可以把他厌恶的某个题目上的讨论卡住;大不了等下一次大会韦伊不在时,再平心静气地把那题目捡起来,甚至第二天他不再作梗时就捡起来。决定由在场的成员一致作出,同时考虑到:一个人与所有其他人意见相左而站在正确一边,绝非不可能(甚至也不罕见)。这条原则对集体工作来说,似乎荒谬。了不起的是:它居然行得通!} (14)。有两次,我相信,我的敏感为此不快,我就走掉了——这是我所知道的、我自己造成的唯一“冲突”迹象。渐渐地,塞尔对团体的影响力变得与韦伊相当。在我还是布尔巴基成员的时候,这并没有造成两人之间的对抗局面,我也没有听说他们后来结下了仇怨。又过了二十五年再回头看,我所认识的五十年代的布尔巴基,在我看来始终是一个范例:在一个围绕共同计划形成的团体里,关系质量上取得了不起的成功。这个团体的品质,在我看来,比它产出的那些书的品质还要罕见。我这一生满是特权,其中之一便是遇到了布尔巴基,并在里面待了几年。我之所以没有留下来,绝不是因为冲突,也不是因为我所说的那种品质败坏了,而是因为一些更私人的工作更强烈地吸引着我,我把全部精力都献给了它们。何况这次离开,既没有给我与团体的关系投下阴影,也没有给我与其中任何一位成员的关系投下阴影。

我该逐一回顾一下我所卷入的冲突情境,那些使我与某位同行或某个学生对立起来的情境——1948到1970年之间。
% ---------------- p.160 ----------------
\origpage{47}\origpage{48}唯一稍稍值得一提的,是与韦伊那两次短暂的闹翻,前面已经说过。还有几片短暂的、极其短暂的阴云,落在我与塞尔的关系上,起因是我的敏感:有时一场谈话不再叫他感兴趣,他便有一种令人有时不知所措的洒脱,一下子把它掐断;或者对我投身其中的某项工作、我坚持的某种看法直说他没兴趣,甚至厌恶——也许我坚持得太多、太频繁了!这从未闹到翻脸的地步。除了气质上的差异,我们的数学意气相投格外强,他必定和我一样感到:我们彼此互补。

与我有着相当甚至更强意气相投的另一位数学家,只有德利涅。说到这儿,我想起1969年德利涅被任命到法国高等科学研究所(IHÉS)一事曾引起紧张,当时我并不觉得那是“冲突”(按说冲突会表现为翻脸,或同行关系中的一个转折)。

我想我已经清点完了——就人与人之间的冲突而言,就可由可见迹象察觉的冲突而言,在我常出入的那个圈子里,同行之间、或同行与学生之间的关系中,二十二年间就只有这些,信不信由你。也就是说,在我选中的那个天堂里没有冲突——那么,想必也就没有轻蔑?数学里又一个矛盾?

看来,我非得再仔细看看不可!

\section{(21) 一个守得很好的公开秘密}

昨天我肯定漏掉了一些小事,比如我与某位同行的关系里那些短暂的“冷淡”,尤其因我的敏感而起。我还该补上三四次我的自尊受挫的场合:有时同行和朋友在他们某篇文章里没有提到,我向他们说过的某个想法或结果想必在他们的工作中起过作用(在我看来如此)。我至今还记得,这说明那是个敏感点,也许并没有随年龄完全消失!除了一次,我都没有向当事人提起,他们的诚意当然不容怀疑。相反的情形肯定也发生过,只是我没听到回声。在我的“小宇宙”里,我不曾得知有哪一桩优先权问题成了翻脸或仇怨的由头,甚至不曾有过当事人之间尖酸刻薄的话。不过,唯一一次我进行了这样的讨论(在我看来那是个明摆着的例子),有过一场口角,倒把气氛澄清了,没有留下怨恨的残渣。那是一位格外出色的同行,他的本领之一,是能以惊人的速度吸收他所听到的一切;我似乎觉得,他常常有一种糟糕的倾向:把刚从别人嘴里听来的想法当成自己的。

这里有一个难题,它在所有数学家身上(而且不止在他们身上)都会以或强或弱的形式出现;它也不只是由于那种自我中心的习气——它促使我们大多数人(我也不例外)去归给自己一些“功劳”,无论真的还是想当然的。对一种情境(数学的或别的)的理解,无论我们怎样达到它,有别人帮助与否,本身都是一件本质上属个人的事,一种个人的经验,其果实是一种眼光,那也必然是个人的。一种眼光有时可以传达,但被传达的眼光不同于原初的眼光。话虽如此,要想辨明他人在自己眼光的形成中占的那一份,却需要极大的警觉。我自己肯定并不总有这份警觉——那是我最不关心的事,然而我却期待别人对我有这份警觉\fnoteP{3}{}{(9月30日)不过,事情的另一个侧面,见6月1日的注(晚于本文三个月),“暧昧”(n° 63),那里考察了对自己和他人有所宽纵的种种陷阱。}!
% ---------------- p.161 ----------------
\origpage{49}迈克·阿廷(Mike Artin)是第一个、也是唯一一个让我听到这件事的人。那天他带着一副泄露公开秘密的玩笑神气对我说:当你终于一把抓住一种实质、并在其中弄懂了点什么的时候,还费劲去分辨哪一份“属于自己”、哪一份“属于别人”,这既不可能,也完全徒劳。这话让我有点不知所措,因为这完全不合乎嘉当、迪厄多内、施瓦茨等人以身教给我的职业伦理。然而我隐隐感到,他的话里、以及他那笑盈盈的目光里,有一个我直到那个年纪都没能领会的真理:我的数学(尤其是数学产出)被自我强烈地占据着,而在迈克身上并非如此。他真给人一种印象:做数学就像一个玩得开心的孩子,同时又不忘吃喝。

\section{(22) 布尔巴基,或我的大运气——及其反面}

在往可见表面之下再潜一点之前,有一个判断现在就已摆在我面前:我五十、六十年代常出入的那个数学圈子,确实是一个“没有冲突的世界”,一点不假!这本身就是相当了不起的事,值得我稍作停留。

我得马上说明:这是一个非常狭小的圈子,是我数学小宇宙的核心部分,局限于我近身的环境中——我定期见面、关系最紧密的那二十来位同行和朋友。逐一检视他们时,我吃惊地发现:其中半数以上是布尔巴基的活跃成员。显然,这个小宇宙的核心与灵魂是布尔巴基。它差不多就是布尔巴基,加上与布尔巴基最亲近的那些数学家。六十年代我自己已不再是这个团体的成员,但我与某些成员的关系仍一如既往地紧密,尤其是与迪厄多内、塞尔、泰特(John Tate)、朗(Serge Lang)、卡蒂耶(Pierre Cartier)。此外我仍常去布尔巴基讨论班,或者不如说,正是那时我才成了常客;也正是在那时,我在那里做了我的大部分报告(关于概形理论)。

大概正是在六十年代,布尔巴基团体里的“调子”滑向了越来越明显的精英主义;我当时肯定是其中一分子,因此也就不可能察觉。我还记得1970年的惊讶:我发现布尔巴基这个名字本身,在数学界广大的阶层(此前为我所不知的阶层)里变得多么不受欢迎——它或多或少成了这些词的同义语:精英主义、狭隘的教条主义、对“典范”形式的崇拜而牺牲活生生的理解、晦涩、阉割式的反自发性,等等等等!而且布尔巴基不只是在“沼泽”里名声不好:六十年代,也许更早,我就偶尔从另一些心性不同、对“布尔巴基风格”过敏的数学家那里听到过回声\fnoteP{4}{(15)}{我并没有得到这样的印象:这种对布尔巴基“风格”的“过敏”,曾给这些数学家与我、或与布尔巴基其他成员及同情者之间的交流造成困难——假如这个团体的精神是宗派精神、是精英中的精英的精神,那倒会如此。在风格与时尚之外,这个团体的所有成员对数学实质都有敏锐的感觉,无论那实质来自何处。只是到了六十年代,我才记得我某位朋友把某些数学家称为“讨厌鬼”,因为他对他们的工作不感兴趣。这类事我几乎一无所知,便倾向于把这样的评价当作真话照收,被那份洒脱的自信镇住了——直到有一天我发现,那位“讨厌鬼”是个独到而深刻的心灵,只是没有福气讨我那位才华横溢的朋友喜欢。我似乎觉得,在布尔巴基某些成员身上,面对别人的工作(当自己不懂或懂得不透时)所持的那种谦逊(至少是保留)态度,已经磨损了。起初,那时那种“数学本能”还残留着——它使人感到一种丰富的实质、一项扎实的工作,而无须参照什么名声或名望。从各处传到我这里的回声看,我似乎觉得,谦逊也好,本能也好,今天在我曾属于的那个数学圈子里都已成了稀罕之物。} (15)。作为一个无条件的信徒,我对此感到惊讶, 
% ---------------- p.162 ----------------
\origpage{50}\origpage{51}也有点难过——我原以为数学能使人心一致!然而我本该记得:我起步的时候,硬吞一篇布尔巴基的文字并不总是轻松、也不总是启发人的,尽管它很干脆。那典范的文字,至少可以这么说,几乎不让人感到它写作时的气氛。如今我似乎觉得,这恰恰是布尔巴基文字的主要缺陷——连偶尔的一丝微笑也不能让人猜到:这些文字是人写的,而且是被远比某种对无情严格典范的无条件效忠更重要的东西连在一起的人写的。

不过,滑向精英主义的问题,以及布尔巴基写作风格的问题,在这里都是题外话。在这里让我吃惊的是:我选作职业圈子的这个“布尔巴基小宇宙”,是一个没有冲突的世界。这事在我看来更了不起,因为这个圈子里的主角们各有强大的数学个性,其中不少被视为“大数学家”,他们每一个都当然有分量形成自己的小宇宙,并且在其中当无可争议的中心和首领!\fnoteP{5}{(16)}{说实话,布尔巴基的成员里有好几位,肯定各有自己那个或大或小的“小世界”,在布尔巴基的小世界之外或之上。但在我自己的情形里,这样一个世界只是在我不再是布尔巴基的一员之后,在我全部精力都投入到属于我个人的工作之后,才在我周围形成——这也许并非偶然。} (16)这些强个性的人,两个十年里在同一个微小宇宙、同一个工作团体里和睦甚至亲切地共处——这才是让我觉得如此了不起、也许是独一无二的事。这与昨天谈布尔巴基时已经浮现的“非凡成功”的印象相合。

看来归根到底,我是有了这份非凡的运气:初次接触数学界时,我恰好落在时空里那个优越的地点——那里几年前刚形成了一个品质非凡的数学圈子,也许就品质而言独一无二。这个圈子成了我的圈子,并始终是我心目中理想“数学共同体”的化身;而这样一个共同体,当时(在我这个化身之外)大概并不比数学史上任何别的时刻更存在,也许除了某些同样狭小的团体(比如也许那个围绕毕达哥拉斯形成、精神却完全不同的团体)。

我对这个圈子的认同非常强,与我四十年代末诞生的数学家新身份不可分割。这是我超出家庭之外第一个被温暖接纳、被当作自己人接受的团体。还有另一条纽带,性质不同:我自己对数学的进路,在这个团体的进路里、在我新圈子里成员的进路里得到了确认。它并不等同于“布尔巴基式”的进路,但显然两者是兄弟。

此外,这个圈子对我而言还该代表那个理想的地点(或差不多是!),那个没有冲突的地点——对它的追寻无疑曾把我引向数学,引向一切科学中最不让我感到有一丝冲突意图的科学!而我刚才说到我的“非凡运气”,心里也明白这份运气有它的反面。如果说它让我得以发展出种种手段,在长辈已成同辈的圈子里以数学家的身份充分施展,那么它同样是一种受欢迎的手段,用来逃避我自己生活里的冲突,并造成一场长期的精神停滞。

\section{(23) 自深处}

% ---------------- p.163 ----------------
\origpage{52}这个“布尔巴基式”的圈子,肯定对我的为人、我对世界以及我在世界中位置的看法产生了强烈影响。这里不是试着界定这种影响、以及它如何在我生活中表现的地方。我只想说:我丝毫也不觉得我那些自命不凡的倾向、以及它们那种唯才是举式的合理化,是由我与布尔巴基的接触、由我进入“布尔巴基圈子”所激发的——至少四十年代末和五十年代不是。那些种子早就在我心里播下,在任何别的圈子里都会找到发展的机会。我报告过的“不中用的学生”那桩事,绝不是这个圈子里主导气氛的典型,恰恰相反,我再说一遍;它只是我自己身上一种暧昧态度的典型。布尔巴基的气氛是一种尊重人的气氛,一种自由的气氛——至少我是这样感受的;这种气氛本性上会打消并减弱一切趋向支配或自命不凡态度的倾向,无论是个人的还是集体的。

这个品质非凡的圈子不复存在了。它什么时候死的,我说不上来;想必没有一个人察觉,也没有一个人在内心深处为它敲响丧钟。我猜想,人身上一定发生了不知不觉的败坏——我们想必都“添了些年头”,变得干硬了。我们成了重要人物,有人听,有权势,有人怕,有人求。火花也许还在,但天真在路上丢失了。我们中也许有人会在死前重新找回它,像一次新生——但那个曾接纳我的圈子已不复存在,我若指望它复活,那是徒劳。一切都回到了秩序之中。

而尊重,也许也在路上丢失了。轮到我们有了学生时,也许已经太晚,最好的东西传不下去了——还有一星火花,但没有了天真,也没有了尊重,除了对“同辈”和“自己人”。

风可以刮起来,吹,烧——我们躲在厚厚的城墙后面,各人带着“自己人”。

一切都回到了秩序之中……

\section{(24) 我的告别,或:外国人}

% ---------------- p.164 ----------------
这场对我数学生涯的回顾,走上了与我预想完全不同的路。说实话,我本来根本没想到要回顾,只想用几行、顶多一两页说说:今天我与那个我已离开的世界是什么关系;也许还反过来说说:从我偶尔收到的回声看,我的旧友与我是什么关系。相反,我原本打算稍仔细地考察一下,我在那些紧张的数学工作年月里引入的某些想法和概念的遭际——它们有时颇为奇怪;或者我该说:我有幸瞥见、并把它们从全然未知的黑夜中拉进半明半暗、有时甚至拉到最明亮的日光之下的那些新类型的对象与结构!如今这番话,在这场已经变成对一段过去的沉思、变成一种更好理解并承担某个有时令人不知所措的当下的努力里,似乎显得走调。看来,原计划中对某个几何“学派”的反思,只好等一个更合宜的时机——那个学派在我的推动下形成,又几乎不留痕迹地挥发掉了\fnoteP{6}{}{这个“更合宜的时机”来得比预想的早,那场反思构成了《收获与播种》第二部分“葬礼”的内容。}。所以眼下,我关心的是把这场回顾进行到底:回顾\origpage{53}\origpage{54}我在数学界里的数学生涯,而不是对一部作品及其命运发议论。

在刚过去的五天里,我被反思笔记之外的别的任务占着,一个回忆却相当执拗地回来了。它将作为我停下的那篇《自深处》的尾声。

事情发生在1977年底左右。几周前,我被传唤到蒙彼利埃轻罪法庭,罪名是“无偿收留并供养一名身份不合法的外国人”(也就是说,一名在法国的居留证件不合规定的外国人)。正是这次传唤,让我知道了1945年那部规定外国人在法地位的法令中令人难以置信的一款;这一款禁止任何法国人以任何形式帮助“身份不合法”的外国人。这条法律,甚至在希特勒德国对犹太人也没有类似的条文,显然从未按字面执行过。由于一个非常奇怪的“巧合”,我有幸被选作第一个试验品,用来首次实施这条独一无二的条款。

有几天我一直发懵,像瘫痪了一样,深深地灰心。突然间,我仿佛被抛回三十五年前——那时人命不值几个钱,尤其是外国人的命……后来我反应过来了,我振作起来。几个月里,我把全部精力投进去,试图动员舆论,先是在我所在的大学和蒙彼利埃,然后在全法范围。正是在这段紧张活动的时期——为了一件后来证明注定失败的事业——发生了我今天可以称之为“我的告别”的那个插曲。

为了在全国范围行动,我给科学界五位特别知名的“人物”(其中一位是数学家)写了信,把这条法律告诉他们——它今天在我看来仍像我被告发那天一样难以置信。我在信里提议共同行动,表明我们反对一条邪恶的法律;这条法律等于把几十万住在法国的外国人置于法律之外,并把几百万其他外国人当作麻风病人一样,让民众对他们心存戒备——这些人一下子成了嫌疑犯,可能给那些不提高警惕的法国人招来最糟糕的麻烦。

奇怪的是,完全出乎我意料:这五位“人物”中没有一位给我回音。看来,我是有些东西要学的……

于是我决定去巴黎,趁布尔巴基讨论班的机会——在那里我准能遇到许多旧友——首先在我最熟悉的数学界里动员舆论。我原以为,这个圈子对外国人的事业会格外敏感,因为我所有的数学同行,和我一样,每天都要与外国同行、学生和留学生打交道;他们大多数、或者全部,都曾有过居留证件上的难处,都曾在警察总局的走廊和办公室领教过专横、而且往往是轻蔑。我曾把计划告诉洛朗·施瓦茨,他说讨论班第一天报告结束后会让我发言,把情况提交给在场的同行。

那天我就这样去了,公文包里一大叠传单,是给同行们的。阿兰·拉斯库(Alain Lascoux)帮我分发,在亨利·庞加莱研究所的走廊里,在第一场之前,以及两场报告之间的“幕间”。如果我没记错,他自己还做了一份小传单——他属于那两三位同行之列:听说了这件事之
% ---------------- p.165 ----------------
\origpage{55}\origpage{56}后,也动了心,在我去巴黎之前就联系了我,表示愿意帮忙\fnoteP{7}{(17)}{我投身的那场行动,得到的热烈回响和切实帮助,主要来自科学界之外。除了阿兰·拉斯库与罗杰·戈德芒的友好支持,我在这里还要特别提到让·迪厄多内的支持:他专程到蒙彼利埃出庭,为这件注定失败的事业在别的证词之外添上他热忱的证词。} (17)。罗杰·戈德芒也在其中,他甚至做了一份传单,标题是“诺贝尔奖得主进监狱?”。这在他是一番好意,可我们显然不在同一个频道上:好像丑闻在于动了一位“诺贝尔奖得主”,而不是随便哪个小人物!

布尔巴基讨论班第一天,人确实很多,极多或多或少与我相识的人,包括布尔巴基的旧友和旧伙伴;我想,他们大多数人想必都在场。还有好几位我过去的学生。我差不多有十年没见过这些人了,来的时候我很高兴有这个机会重见他们,哪怕一下子见这么多!不过,总会有机会再小聚的……

然而这次重逢“不是那么回事”,从一开始就相当清楚。许多手伸过来、握上了,这是肯定的;许多“哟,你也在这儿,什么风把你吹来了?”的问题,也是有的——但在快活的语调背后,有一种说不清道不明的尴尬气息。是因为我带来的事业他们并不真正感兴趣,而他们来这里是为了一场一年三次的数学典礼,需要全神贯注?还是与我来做什么无关,而是我这个人本身引起了那种尴尬——有点像一位还俗的神父出现在一群纯正的神学院学生中间所引起的那种尴尬?我说不好——也许两者兼有。在我这一边,我忍不住注意到,某些曾经熟悉、甚至曾是我的朋友的面孔上发生的变化。他们仿佛僵住了,或者塌陷了。我曾在他们脸上认识的那种灵动似乎消失了,仿佛从不曾有过。我觉得自己像是面对一些陌生人,仿佛从来没有什么把我与他们连在一起。我隐约感到,我们并不生活在同一个世界里。我原以为借这个不寻常的机会能找回一些兄弟,结果面对的却是陌生人。教养很好,这得承认:我不记得有过尖酸的话,也不记得有传单被丢在地上。事实上,好奇心帮了忙,分发的传单(几乎)都该被读过了。

然而那条邪恶的法律并没有因此受到威胁!我得到了我的五分钟,也许还占了十分钟,来谈那些在我看来是兄弟、却被叫作“外国人”的人的处境。阶梯教室里坐满了同行,比我做数学报告时还安静。也许对他们讲话的那份确信,已经不在我身上了。不再有从前那种同情与兴趣的流动。人群中想必有些赶时间的人,我一定是这样对自己说的;我缩短了发言,提议立刻与感到关切的同行们留下来,更详细地商议可以做些什么。

宣布散会时,是全体向出口的冲锋——显然,人人都有一班火车或地铁马上要开,绝不能错过!一两分钟之内,埃尔米特阶梯教室就空了,简直不可思议!空荡荡的大阶梯教室里,刺眼的灯光下,只剩下三个人,包括阿兰和我。第三个人我不认识,我敢打赌又是那种见不得人的外国人,交往可疑,而且身份还不合法!我们没有花时间对刚才在眼前上演的那一幕说得太多——它本身已经够雄辩了。也许我还是唯一一个不敢相信自己眼睛的人,而我的两位朋友当时很得体地对此不加评论。显然,我是初来乍到……

晚上在阿兰和他前妻雅克琳家里结束,盘点一下形势,把可以做的事过一遍;也彼此多认识一点。无论是这一天还是后来,我都没有费心把刚经历的这个插曲放进过去的脉络里去看。然而正是那一天,我一定是不言不语地明白了:某个圈子、某个我曾认识并爱过的世界已经不存在了;我原以为能重新找回的一份活生生的温暖,早已消散——多半消散很久了。

% ---------------- p.166 ----------------
这并没有阻止那个世界一年一年传到我这里的回声——那个温暖已逃逸的世界——许多次让我不知所措,让我痛苦地被打动。我怀疑这场反思对将来能否对此有所改变——除非也许是:我再被这样触动时,会少梗着脖子抗拒一点……
